\hypertarget{oppgave3}{}
\section*{Oppgave 3}
\addcontentsline{toc}{section}{Oppgave 3}

Tidsfunksjonene (a)--(e) kobles sammen med frekvensspektrene (1)--(5) ved å sammenligne DC-verdi, frekvens, periodetid og fase.

% --------------------------------------------------
\subsection*{a)}
Tidsfunksjon (a) har en DC-komponent på omtrent 2.
Den må derfor tilhøre spektrum (1) eller (3), siden begge disse har en DC-komponent på 2.
Spektrum (3) har frekvenskomponenter ved:
\[
f=\pm1.2\text{ Hz}.
\]
Komponentene har magnitude 1.5, slik at cosinusleddet får amplituden:
\[
A=2\cdot1.5=3.
\]
Fasen ved positiv frekvens er $-\pi/4$, og signalet kan derfor skrives som:
\[
x(t)=2+3\cos\left(2\pi1.2t-\frac{\pi}{4}\right).
\]
Ved $t=0$ blir signalverdien:
\[
x(0)=2+3\cos\left(-\frac{\pi}{4}\right)\approx4.12.
\]
Dette stemmer med tidsfunksjon (a).
Dermed kobles:
\[
\boxed{\text{(a)}\rightarrow\text{(3)}}.
\]
% --------------------------------------------------
\subsection*{b)}
Tidsfunksjon (b) har ingen DC-komponent og består av én tydelig frekvens.
Spektrum (5) inneholder kun frekvenskomponenter ved:
\[
f=\pm1.5\text{ Hz}.
\]
Komponentene har magnitude 1.5, slik at amplituden til cosinusleddet blir:
\[
A=2\cdot1.5=3.
\]
Fasen ved positiv frekvens er $\pi$, og signalet kan derfor skrives som:
\[
x(t)=3\cos(2\pi1.5t+\pi).
\]
Ved $t=0$ blir signalverdien:
\[
x(0)=3\cos(\pi)=-3.
\]
Dette stemmer med tidsfunksjon (b).
Dermed kobles:
\[
\boxed{\text{(b)}\rightarrow\text{(5)}}.
\]

% --------------------------------------------------
\subsection*{c)}
Tidsfunksjon (c) har en DC-komponent på omtrent 2.
Dette passer med spektrum (1), som har en DC-komponent på 2 og frekvenskomponenter ved:
\[
f=\pm1.2\text{ Hz}.
\]
Komponentene har magnitude 1.5, slik at cosinusleddet får amplituden:
\[
A=2\cdot1.5=3.
\]
Fasen ved positiv frekvens er $\pi/2$, og signalet kan derfor skrives som:
\[
x(t)=2+3\cos\left(2\pi1.2t+\frac{\pi}{2}\right).
\]
Ved $t=0$ blir signalverdien:
\[
x(0)=2+3\cos\left(\frac{\pi}{2}\right)=2.
\]
Dette stemmer med tidsfunksjon (c).
Dermed kobles:
\[
\boxed{\text{(c)}\rightarrow\text{(1)}}.
\]
% --------------------------------------------------
\subsection*{d)}
Spektrum (2) inneholder frekvenskomponenter ved 0.6 Hz og 1.5 Hz.
Begge frekvensene er heltallsmultipler av en felles grunnfrekvens:
\[
0.6=2\cdot0.3,\qquad1.5=5\cdot0.3.
\]
Grunnfrekvensen blir derfor:
\[
f_0=0.3\text{ Hz}.
\]
Periodetiden blir:
\[
T_0=\frac{1}{f_0}
=\frac{1}{0.3}
\approx3.33\text{ s}.
\]
Dette stemmer med periodetiden til mønsteret i tidsfunksjon (d).
Dermed kobles:
\[
\boxed{\text{(d)}\rightarrow\text{(2)}}.
\]
% --------------------------------------------------
\subsection*{e)}
Spektrum (4) inneholder frekvenskomponenter ved 1.2 Hz og 2 Hz.
Begge frekvensene er heltallsmultipler av en felles grunnfrekvens:
\[
1.2=3\cdot0.4,\qquad2=5\cdot0.4.
\]
Grunnfrekvensen blir derfor:
\[
f_0=0.4\text{ Hz}.
\]
Periodetiden blir:
\[
T_0=\frac{1}{f_0}
=\frac{1}{0.4}
=2.5\text{ s}.
\]
Dette stemmer med periodetiden til mønsteret i tidsfunksjon (e).
Dermed kobles:
\[
\boxed{\text{(e)}\rightarrow\text{(4)}}.
\]
Samlet blir koblingen:
\[
\boxed{
\text{(a)}\rightarrow\text{(3)},\qquad
\text{(b)}\rightarrow\text{(5)},\qquad
\text{(c)}\rightarrow\text{(1)},\qquad
\text{(d)}\rightarrow\text{(2)},\qquad
\text{(e)}\rightarrow\text{(4)}
}
\]